\documentclass[11pt,a4paper]{article}

\usepackage[utf8]{inputenc}
\usepackage[T1]{fontenc}
\usepackage{amsmath,amssymb,amsthm}
\usepackage{booktabs}
\usepackage{array}
\usepackage{geometry}
\usepackage{hyperref}
\usepackage{xcolor}
\geometry{margin=2.5cm}

\newcommand{\sbx}{\textsf{S-box}}
\newcommand{\sbxes}{\textsf{S-boxes}}
\newtheorem{remark}{Remark}

\title{\textbf{Differential Active-S-box Analysis of the Krakken-2048 SPN Core:\\
A Verified Two-Round Lower Bound and Its Limitations}}
\author{Working note}
\date{\today}

\begin{document}
\maketitle

\begin{abstract}
This note records the results of an automated active-S-box analysis of
Krakken-2048, carried out with a byte-lane Mixed-Integer Linear Programming
(MILP) model, covering both \emph{differential} and \emph{linear} trails and
both the SPN+XRBD subset and the complete round function. For the full
(SPN+XRBD) round we prove the minimum differential active-S-box count at every
round from one to eight: $5, 37, 69, 101, 133, 165, 197, 229$. Each figure was
obtained with the branch-and-bound optimality gap closed to zero (a proof, not a
time-limited estimate), and the two-round SPN-core trail was independently
verified layer-by-layer. The complete-round model (additionally including the
PRESSURE/ARX and InkCloud layers) and the linear-trail model are proven to give
the same counts at every round from one to eight---the expected behaviour for a
wide-trail design with a maximal MDS branch number, where differential and
linear active-S-box counts coincide and the post-Chi mixing layers cannot lower
the floor set by XRBD.

The headline consequence: the full eight-round permutation has \textbf{no
differential characteristic with fewer than $229$ active S-boxes}, bounding any
eight-round differential characteristic probability by $2^{-1374}$ and thereby
ruling out single-characteristic differential cryptanalysis with very large
margin; the linear bound, where confirmed, gives the corresponding correlation
bound $2^{-1374}$ via the S-box maximum squared correlation of $2^{-6}$. From two
rounds onward each additional round adds exactly $32$ active S-boxes. We remain
explicit about scope: the per-S-box bounds use the standard active-S-box
argument (not a cluster/hull analysis), and the PRESSURE layer is modelled
conservatively (modular addition as activity-spreading, ignoring carries), which
keeps all reported counts sound lower bounds. We also record a methodological
finding: the open-source solver SCIP closes these instances where CBC's lower
bound stays frozen, making proven bounds attainable without a commercial solver.
\end{abstract}

\section{Scope and methodology}

We model differential propagation through Krakken-2048 at \emph{byte-lane}
granularity: the $2048$-bit state is $32$ words $\times\,8$ byte-lanes
$=256$ lanes, one binary activity variable per lane (equivalently, one per
\sbx). A MILP minimises the total number of active \sbxes\ over $N$ rounds
subject to each layer's activity-propagation constraints. When the solver
reports an optimum \emph{with a closed branch-and-bound gap}, the objective is a
proven lower bound on the number of active \sbxes\ for any nonzero differential
characteristic over those $N$ rounds.

\paragraph{Why byte-lane and not word level.}
An earlier word-level model (one variable per $64$-bit word) yielded a flat
$4$ active words per round, because at word granularity the bit-rotations
($\rho$, the XRBD rotations, the ARX shifts) are invisible: a rotation maps an
active word to an active word and nothing more. Byte-lane granularity captures
the essential effect that a rotation by $r\not\equiv 0 \pmod 8$ spreads an active
lane across two adjacent output lanes, which is what makes activity compound
across rounds.

\paragraph{Conservative (sound) modelling.}
Every modelling choice under-counts activity, so the resulting bound is a valid
\emph{lower} bound:
\begin{itemize}\itemsep2pt
  \item Rotations are modelled as activity \emph{spread} (an output lane is
        active if any input lane rotating into it is active).
  \item The MDS layer is modelled by its branch-number constraint: for each row
        and lane, the count of active input bytes plus active output bytes is
        either $0$ or at least $9$.
  \item $\Theta$ is modelled as non-cancelling column-parity spread.
  \item $\chi$ is modelled with its serial cross-coupling and the $32$-bit
        half-swap.
\end{itemize}

\paragraph{MDS provenance.}
The $8\times 8$ circulant MDS uses the Whirlpool coefficient row
$(1,1,4,1,8,5,2,9)$ over $\mathbb{F}_{2^8}$ with reduction polynomial
\texttt{0x11D}, which has a proven branch number of $9$ (the maximum for an
$8\times 8$ MDS matrix). We rely on this externally-established property rather
than re-deriving it; we do, however, verify below that the matrix \emph{as
implemented} enforces branch~$9$ inside the round.

\section{Result: a verified two-round bound}

\begin{table}[h]
\centering
\begin{tabular}{lcccc}
  \toprule
   & \textbf{Diff.\ SPN+XRBD} & \textbf{Diff.\ all layers} & \textbf{Linear all layers} & \\
  \textbf{Rounds} & \textbf{(active \sbxes)} & \textbf{(active \sbxes)} & \textbf{(active \sbxes)} & \textbf{DP / corr$^2$ bound} \\
  \midrule
  1 & 5   & 5   & 5   & $2^{-30}$   \\
  2 & 37  & 37  & 37  & $2^{-222}$  \\
  3 & 69  & 69  & 69  & $2^{-414}$  \\
  4 & 101 & 101 & 101 & $2^{-606}$  \\
  5 & 133 & 133 & 133 & $2^{-798}$  \\
  6 & 165 & 165 & 165 & $2^{-990}$  \\
  7 & 197 & 197 & 197 & $2^{-1182}$ \\
  8 & \textbf{229} & \textbf{229} & \textbf{229} & $\mathbf{2^{-1374}}$ \\
  \bottomrule
\end{tabular}
\caption{Minimum active-\sbx\ counts, all entries proven to optimality (solver
gap closed to $0$). ``Diff.\ SPN+XRBD'' is the differential bound for
Theta$\to$MDS$\to$Rho$\to$Pi$\to$Chi$\to$XRBD. ``Diff.\ all layers'' and
``Linear all layers'' additionally include the PRESSURE (ARX) and InkCloud
layers. All three columns coincide at every round---the expected behaviour for a
wide-trail design with a maximal MDS branch number: differential and linear
active-S-box counts agree, and the post-Chi mixing layers do not lower the count
set by XRBD. The bound column is $(2^{-6})^{\text{count}}$: the
differential-characteristic probability bound and, by wide-trail symmetry, the
linear correlation-squared bound. For reference the SPN core alone (no XRBD)
gives $5,27,55$ at rounds one to three.}
\label{tab:bounds}
\end{table}

\paragraph{Headline result.}
The eight-round count of $229$ active S-boxes proves that no eight-round
differential characteristic has probability above $2^{-1374}$. This is proven for
the complete round function (all layers) and, by the linear model, the same
bound holds for linear cryptanalysis (correlation-squared $\le 2^{-1374}$). For
an eight-round characteristic to threaten a $2048$-bit-state construction its
probability would need to be far higher; the proven bound is below this by a
margin of more than a thousand bits, so single-characteristic differential and
linear cryptanalysis are both ruled out conclusively. These are direct proofs at
eight rounds, not extrapolations.

\paragraph{Differential/linear coincidence and the all-layers model.}
At every round from one to eight, the complete-round differential bound and the
linear bound equal the SPN+XRBD differential bound (proven to optimality). This
is the expected behaviour of a wide-trail design built on a maximal-branch ($9$)
MDS matrix: at the byte-lane activity level the linear model is the transpose of
the differential model, and the MDS branch number is identical for a matrix and
its transpose, so the two active-S-box counts coincide; and the PRESSURE/InkCloud
layers, being additional mixing, cannot lower the active-S-box floor that XRBD
already forces. We therefore present the coincidence as a \emph{confirmation} of
wide-trail behaviour rather than as an independent rediscovery: the linear and
complete-round models share the activity-level structure of the differential
model, viewed from a different direction.

\paragraph{The linear trend and the XRBD contribution.}
From two rounds onward each additional round adds \emph{exactly} $32$ active
S-boxes. At one round the SPN-core and full-round counts coincide ($5$), as
expected---the S-boxes are counted at the Chi input, before the post-Chi layers
act. Beyond one round the full-round count exceeds the SPN core by a growing
margin ($+10$ at two rounds, $+14$ at three), confirming that XRBD's cross-word
rotations break the cheap lane-aligned trail the SPN core admits (a trail that
rides the MDS at its exact branch number along a single byte-lane; see the Remark
below). The flat $+32$ increment over eight proven data points is a strong
empirical regularity, but we emphasise it is observed, not derived: the
eight-round value rests on a direct proof, and we make no theoretical claim that
the increment would remain $32$ beyond eight rounds.

\subsection{Independent verification of the trail}

Because the bound is only as trustworthy as the encoding, the optimal trail was
extracted and re-checked outside the MILP. Each layer was re-simulated in
independent code on the solved input pattern and compared against the solver's
layer variables; all layers matched. The branch-number property was checked on
every active MDS row-lane.

\paragraph{Minimal trail structure.}
The minimum-weight input difference has $4$ active bytes, all in the same
byte-lane (lane~$7$), located in words $0, 12, 19, 28$. Its propagation is:
\begin{itemize}\itemsep2pt
  \item \textbf{Round 0:} the $4$ active bytes lie in $4$ distinct MDS rows; each
        active row realises the branch number \emph{exactly} ($8$ active inputs,
        $1$ active output, sum $9$), yielding $9$ active \sbxes.
  \item \textbf{Round 1:} activity spreads to $16$ active MDS row-lanes, again
        each at exactly branch $9$ ($8$-in, $1$-out), yielding $18$ active
        \sbxes.
\end{itemize}
The total is $9 + 18 = 27$, matching the solver objective. Every active MDS
row-lane satisfied $\text{in}+\text{out}\ge 9$ with zero violations, confirming
the Whirlpool branch number is enforced by the implementation.

\begin{remark}
The minimal trail rides the MDS at its \emph{exact} branch number throughout and
remains aligned to a single byte-lane. The per-round MDS activity grows as
$9 \to 18$, i.e.\ linearly (factor $2$) at the cheapest corner of the diffusion,
rather than super-linearly. This is precisely the class of structured,
position-aligned trail that the XRBD butterfly layer is intended to disrupt.
\end{remark}

\section{What this does \emph{not} establish}
\label{sec:limits}

We state the limitations plainly, as each bears directly on whether the 8-round
parameter is justified.

\begin{enumerate}\itemsep3pt
  \item \textbf{Linear bound is structurally coincident, not independent.} The
        linear-trail MILP (the transpose of the differential model) is proven at
        all eight rounds and reproduces the differential count exactly, as
        wide-trail theory predicts. Because the activity-level linear model shares
        the differential model's structure, it \emph{confirms} rather than
        independently re-derives the bound; a genuinely independent check would
        require a bit-level mask model with explicit transposed linear maps
        (future work).

  \item \textbf{Single characteristic, not differentials/hulls.} The argument
        bounds the best single trail. Differential attacks can in principle
        exploit clusters of characteristics, and linear attacks linear hulls. The
        $2^{-1374}$ margin leaves so much headroom that cluster/hull effects are
        extremely unlikely to bridge it. An empirical differential-clustering test
        on the minimal-trail input difference (companion experiments note,
        $200{,}000$ pairs per round, $128$-bit output projection) found no
        detectable cluster at any round from one to eight (detection floor
        $\approx 2^{-17}$), giving direct evidence against a fat cluster on this
        trail; a full analytical cluster/hull quantification across all
        input differences remains open.

  \item \textbf{PRESSURE modelled conservatively.} The complete-round model
        includes the PRESSURE (ARX) layer with modular addition modelled as
        activity-spreading (carry ignored). This can only \emph{add} activity, so
        the reported counts remain sound lower bounds; a carry-aware ARX model
        could only raise them.

  \item \textbf{The flat $+32$ increment is observed, not proven in general.}
        Every value through eight rounds is a direct proof, so the eight-round
        figure does not depend on the trend. We make no claim that the per-round
        increment remains $32$ beyond eight rounds.
\end{enumerate}

\paragraph{Net assessment.}
On the differential metric the eight-round parameter is directly justified:
$229$ active S-boxes and a $2^{-1374}$ characteristic bound place the design far
beyond the reach of differential cryptanalysis. The complete-round model (all
layers) and the linear model are proven to give the same counts at every round
from one to eight, extending the conclusion to the full round function and to
linear cryptanalysis under the wide-trail coincidence described above. The
remaining refinements---a bit-level independent linear model, a carry-aware ARX
model, and cluster/hull estimates---are expected to preserve or strengthen the
picture, not weaken it.

\section{Methodological note: solver choice is decisive}
\label{sec:method-note}

The bounds above were not obtainable with the default open-source MILP solver.
With CBC, the lower bound (``best possible'') stayed frozen at the LP-relaxation
value $73/512 \approx 0.143$ for the entire run at three rounds: the solver found
feasible trails (upper bounds $57 \to 56 \to 55$) but generated no effective
cutting planes, so the gap never closed and the run terminated on its time limit
with a meaningless lower bound. The two-round instance was provable under CBC
only because its branch-and-bound tree was small enough to exhaust by
enumeration.

SCIP, by contrast, generated tens of thousands of cuts and moved the dual bound
steadily ($0.14 \to 8 \to 22 \to 35 \to 55$), closing the three-round gap to $0$
in roughly $51$ minutes on a single core. We also tested a hand-added
convex-hull MDS lifting ($\sum_{\text{row}} v \ge 9\,v_k$ for each variable);
it preserved correctness (two-round result unchanged at $27$) but did
\emph{not} accelerate convergence---it left the root LP relaxation essentially
unchanged because that relaxation's weakness is a uniform-fractional-activity
artifact spanning all lanes, not a per-row deficiency. The practical conclusion
is that a strong cut-generating solver (SCIP, or commercial alternatives) is the
decisive ingredient for these instances; formulation tightening of the per-row
MDS constraint alone is not.

\section{Reduced-round distinguisher experiments}
\label{sec:distinguishers}

The permutation was subjected to two classic structural distinguisher tests
using a reduced-round harness (\texttt{harness.c}) driven by the
\texttt{run.sh} build-and-run wrapper. Both tests operate directly on the
\emph{permutation} (not the sponge) at a configurable round count and use a
deterministic splitmix64 PRNG seeded with the round count for reproducibility.

\paragraph{Test 1: Integral / saturation distinguisher.}
A structure of $2^8 = 256$ inputs is formed by saturating one chosen input
byte over all $256$ values while holding every other byte constant at a random
value. The reduced-round permutation is applied to each element and all $256$
outputs are XOR-summed. For an ideal permutation, each output byte is balanced
(XOR-sum $= 0$) with probability $1/256$ by chance, giving an expected
baseline of $\mathit{structures}/256$ zero-sum events per output byte position
across many random structures. A byte position that is zero far more often
than this baseline---specifically, more than $8\times$ the ideal \emph{and}
at least $4$ hits---is flagged as a possible distinguisher.

\paragraph{Test 2: Zero-sum / cube distinguisher.}
A $d$-dimensional input subspace is formed by choosing $d$ random bit
positions in the 2048-bit state and sweeping all $2^d$ combinations of those
bits (holding the remaining bits constant at a random value). The outputs over
the full cube are XOR-summed. A sum of zero across \emph{all} $256$ output
bytes constitutes a zero-sum distinguisher. Multiple random cube placements
are tried per round and the best (most zero output bytes) is reported.

\paragraph{Results.}
Table~\ref{tab:distinguishers} records all runs. At every round count from
$1$ to $4$, both tests are consistent with the random-permutation baseline.
No integral distinguisher was detected at any round: the highest observed
ratio of balanced bytes to the ideal baseline was $1.8\times$ (rounds 3
and 4), well below the $8\times$ detection threshold. No zero-sum
distinguisher was detected: the best zero-sum result across all trials was
$4/256$ output bytes summing to zero (round 4, cube dimension $18$), against
an ideal expectation of approximately $1$.

\begin{table}[h]
\centering
\renewcommand{\arraystretch}{1.2}
\begin{tabular}{cccccc}
  \toprule
  \textbf{Rounds} &
  \textbf{Structures} &
  \textbf{Integral best ratio} &
  \textbf{Cube dim $d$} &
  \textbf{Cube trials} &
  \textbf{Zero-sum best} \\
  \midrule
  1 & 4096 & $1.6\times$ ideal & 10 &  8 & $3 / 256$ \\
  2 & 4096 & $1.7\times$ ideal & 10 &  8 & $2 / 256$ \\
  3 & 4096 & $1.8\times$ ideal & 14 & 16 & $3 / 256$ \\
  4 & 4096 & $1.8\times$ ideal & 18 & 16 & $4 / 256$ \\
  \bottomrule
\end{tabular}
\caption{Reduced-round integral and zero-sum distinguisher results.
``Integral best ratio'' is the highest observed ratio of balanced-byte
events to the ideal baseline ($\mathit{structures}/256 = 16$) across all
$256$ output byte positions; the flagging threshold is $8\times$.
``Zero-sum best'' is the maximum number of output bytes that summed to zero
across all cube placements in that run; the ideal expectation for a random
permutation is $\approx 1$.  No distinguisher was detected at any round
count.  The cube dimension was increased with the round count to compensate
for the growing algebraic degree: $2^{18}$ inputs per trial at round~4.}
\label{tab:distinguishers}
\end{table}

\paragraph{Interpretation.}
The integral result confirms that no output byte is forced to balance by
the permutation structure at any of the tested round counts: even at one
round, the XOR-sum distribution over saturated structures is consistent with
a random bijection. The zero-sum result is consistent with the algebraic
degree growing rapidly with the round count: a zero-sum distinguisher
requires a cube of dimension at least equal to the algebraic degree of the
output in the chosen input variables, and no such cube was found up to
dimension~$18$ at four rounds ($2^{18} = 262{,}144$ evaluations per trial,
$16$ independent placements). These results complement the MILP active-S-box
bounds: the MILP bounds rule out differential and linear distinguishers
analytically; the harness results give direct empirical evidence against
integral and algebraic zero-sum distinguishers at reduced round counts.

\paragraph{Division-property integral (analytical, MILP).}
The two tests above are empirical samplings. Their analytical counterpart is the
bit-based division property (two-subset CBDP), which proves an output bit balanced
exactly when its division trail through the round function is infeasible---a
proof, not a statistical observation. We have built a sound MILP harness for this:
each round is modelled bit-for-bit (the linear layers by their exact GF(2)
matrices, the ARX PRESSURE layer by a sound full-adder over-approximation, the
ABYSSAL \sbx\ by its division transitions), with the circuits validated bit-for-bit
against the compiled permutation. The practical obstacle was the \sbx\ model: the
standard ``selector'' encoding adds $\approx\!2015$ binary variables per \sbx\
($\approx\!5\times10^{5}$ per round), and at that size the solver could not resolve
a single output bit. Replacing it with a compact, \emph{exact} inequality model
($661$ inequalities, no auxiliary variables, verified over all $2^{16}$ points of
the \sbx\ division polytope) cuts the one-round program from $937{,}952$ to
$422{,}112$ variables and makes it solvable: round~$1$ now solves to optimality
(output bit~$0$ proved reachable in $\approx\!160$\,s) where the selector model
never terminated. The round-$1$ result is consistent with the empirical integral
finding---first-round diffusion already balances away---and the exhaustive
multi-round sweep that would establish the integral-distinguisher length is an
ongoing computation. The harness is released with this analysis for reproduction
and extension.

\section{Rotational and RX-difference distinguisher experiments}
\label{sec:rotational}

The one-round permutation was probed for rotational and RX-difference structure
using a dedicated command-line tool (\texttt{rotational}), run in heavy mode
(\texttt{--heavy --no-timeout}) with 500\,000 trials per configuration and the
timeout disabled.  Four test categories were executed.

\paragraph{Test 1: Word-level rotational distinguisher.}
Pairs of inputs related by a rotation of $r \in \{1,2,3,4\}$ words were fed to
the one-round permutation and the fraction of output pairs preserving the same
rotation was recorded.  For all four rotation offsets, the number of consistent
output words was $0$ against a random baseline of ${\approx}3.47\times10^{-13}$.
No word-level rotational symmetry was detected.

\paragraph{Test 2: Bit-level rotational distinguisher.}
The same procedure was applied at bit granularity for rotation offsets
$r \in \{1, 7, 13, 19, 32\}$ bits.  All five offsets yielded $0$ consistent
output words.  No bit-level rotational symmetry was detected.

\paragraph{Test 3: Single-word RX-difference distinguisher.}
For rotation constants $\mathit{rx\_rot} \in \{7, 13, 17, 19, 31, 37\}$ and
injection words $\{0, 8, 16\}$, the maximum per-output-bit bias across all $2048$
output bits was measured over $500{,}000$ trials.  The $5\sigma$ detection
threshold at this sample size is $0.003536$.  In $17$ of the $18$ configurations
the maximum bias was well below threshold (range $0.002228$--$0.002796$).

One configuration (\texttt{rx\_rot=17}, \texttt{inj\_word=8}) yielded a maximum
bias of $0.003572$ at bit~$1421$, flagging $1$ out of $2048$ bits above the
$5\sigma$ threshold.  Given the number of bits tested across the full suite, a
single marginal exceedance at ${\approx}1.01\times$ the threshold is consistent
with the expected false-positive rate under the null hypothesis and is not
considered a genuine distinguisher candidate; see the heavy-mode sweep below for
independent confirmation.

\paragraph{Test 4: Heavy-mode PRESSURE-block RX sweep.}
A systematic sweep over all $63$ non-trivial rotation constants was conducted for
four multi-word injection configurations: \texttt{blkA}~(words~0--3),
\texttt{blkB}~(words~16--19), \texttt{blkA+B}~(words~0--3 and~16--19), and
\texttt{w0-7}~(words~0--7), each at $500{,}000$ trials per rotation constant
($252$ jobs in total; wall-clock time ${\approx}962\,\mathrm{s}$).  The five
highest-ranking results are shown in Table~\ref{tab:rotational}.  No configuration
exceeded the $5\sigma$ threshold, and the highest bias observed across the entire
sweep was $0.003472$---$1.8\%$ below the $5\sigma$ threshold of $0.003536$.

\begin{table}[h]
\centering
\renewcommand{\arraystretch}{1.2}
\begin{tabular}{clccccc}
  \toprule
  \textbf{Rank} & \textbf{Config} & \textbf{rot} &
  \textbf{max bias} & \textbf{$5\sigma$ thresh} &
  \textbf{Anomalous bits} & \textbf{Top bit} \\
  \midrule
  1 & \texttt{blkB(w16-19)} & 48 & 0.003472 & 0.003536 & 0 & 1873 \\
  2 & \texttt{w0-7}         & 21 & 0.003198 & 0.003536 & 0 & 1024 \\
  3 & \texttt{blkA+B(w0-3,16-19)} & 49 & 0.003162 & 0.003536 & 0 & 1883 \\
  4 & \texttt{w0-7}         & 14 & 0.003146 & 0.003536 & 0 &  973 \\
  5 & \texttt{blkB(w16-19)} & 12 & 0.003140 & 0.003536 & 0 &  442 \\
  \bottomrule
\end{tabular}
\caption{Top-5 results from the heavy-mode RX sweep over all 63 rotation
constants and 4 multi-word injection configurations ($500{,}000$ trials each).
No configuration produced any output bit with bias above the $5\sigma$ threshold
of $0.003536$.  The highest observed bias of $0.003472$ is $1.8\%$ below the
detection threshold.}
\label{tab:rotational}
\end{table}

\paragraph{Interpretation.}
All four test categories return results consistent with the random-permutation
null hypothesis.  The marginal single-bit exceedance observed in the single-word
sweep (Test~3) is not reproduced in any of the $252$ independent configurations
of the heavy sweep (Test~4), confirming it as a statistical fluctuation rather
than a structural anomaly.  At one round, no rotational or RX-difference
distinguisher is detectable at the tested trial counts and significance level.

\section{Empirical cryptanalytic probe suite}
\label{sec:attacks}

The permutation was subjected to a five-test empirical probe suite
(\texttt{attacks.c}), targeting attack families not directly addressed by the
MILP active-S-box analysis: impossible differentials, algebraic degree growth,
boomerang distinguishers via the BCT, differential-linear bias, and invariant
subspace / S-box coset structure.  Tests were run at one and two rounds in
normal mode (timeouts enabled) with the default trial budgets.  Results are
summarised in Table~\ref{tab:attacks}; the subsections below discuss each test
in detail.

\begin{table}[h]
\centering
\renewcommand{\arraystretch}{1.2}
\begin{tabular}{clp{3.5cm}p{4cm}}
  \toprule
  \textbf{Test} & \textbf{Attack family} & \textbf{1 round} & \textbf{2 rounds} \\
  \midrule
  1 & Impossible differential  & Clean            & Clean \\
  2 & Algebraic degree growth  & $\ge 23$         & $\ge 21$ \\
  3 & BCT / boomerang          & BCT max $= 6$\,$^\dagger$; no full-cipher distinguisher & (same) \\
  4 & Differential-linear      & Skipped ($r < 2$) & Clean; max $|$corr$| = 0.0034$ \\
  5 & Invariant subspace       & Clean            & Clean \\
  \bottomrule
\end{tabular}
\caption{Summary of empirical probe results at one and two rounds.
$^\dagger$~The BCT maximum of $6$ is a property of the S-box alone and is
round-independent; it does not translate to a full-permutation distinguisher at
any tested round count.}
\label{tab:attacks}
\end{table}

\subsection{Test 1: Impossible differential}

For each of four input-bit positions ($0, 64, 512, 1984$) and four output-byte
positions ($0, 31, 127, 255$), $200{,}000$ random pairs $(x,\,x \oplus \delta)$
were encrypted and the output-difference histogram at the chosen output byte was
recorded.  Any output difference with a histogram count of zero constitutes an
impossible-differential candidate.

At one round all completed input-bit configurations were clean ($0/255$
non-trivial differences missing); the 30-second timeout interrupted the run
before the final input bit ($\mathit{input\_bit} = 1984$) could be completed.
At two rounds the timeout fired after two of the four input bits, again with no
candidate found.  Neither run produced an impossible-differential candidate.

\subsection{Test 2: Algebraic degree growth}

The multilinear degree of chosen output bits over $\mathrm{GF}(2)$ was measured
using the top-coefficient (higher-order derivative) method.  The test climbs
dimension $d$ from $1$ upward, running $5{,}000$ inner trials per step; each
step at dimension $d$ costs $2^d$ permutation evaluations per trial.  It stops
at the first $d$ where no nonzero top coefficient is found, reporting the last
confirmed $d$ as an empirical lower bound.

At one round, all inner trials returned a nonzero top coefficient at every
$d$ from $1$ to $23$; the run was cut by the 120-second timeout before
completing $d = 24$.  At two rounds the confirmed lower bound is $\ge 21$,
with the timeout firing at $d = 22$ (two-round evaluations are more expensive).

Both results are consistent with the theoretical expectation: the full-state
algebraic degree saturates to approximately $2048$ within one to two rounds,
owing to the S-box's per-byte degree of $7$ combined with the diffusion layers.
The timeouts are not failures; they reflect the exponential cost of reaching
high $d$ with the default 120-second cap rather than any absence of degree.

\subsection{Test 3: S-box BCT and boomerang distinguisher}

\paragraph{Part A: Boomerang Connectivity Table.}
The full $256 \times 256$ BCT of the Abyssal S-box was computed.  The maximum
entry (excluding $(a,b) = (0,\ast)$ and $(\ast,0)$) is $6$, occurring at
$(a, b) = (\texttt{0x01}, \texttt{0x05})$.  The distribution of non-zero
entries is: $6$ at $510$ pairs, $4$ at $255$ pairs, $2$ at $31{,}620$ pairs.

A BCT maximum of $6$ exceeds the DDT-flat ideal of $2$ (probability
$2^{-7}$ per switch), giving a per-S-box boomerang switch probability of
$2^{-5.42}$.  This is a property of the S-box alone and is
round-independent.  It warrants noting but, as Part B shows, does not propagate
to a full-permutation distinguisher.

\paragraph{Part B: Empirical boomerang probe.}
A $50{,}000$-trial empirical boomerang quartet experiment was run at one and two
rounds using input difference $\alpha = \texttt{0x01}$ at byte $0$.  At one
round the maximum output-difference count was $232$ against a uniform expectation
of $195.3$ (ratio $1.19$); at two rounds the count was $229$ (ratio $1.17$).
Both are well below the $2.0$ anomaly threshold.  No full-permutation boomerang
distinguisher was detected at either round count.

\subsection{Test 4: Differential-linear distinguisher}

This test requires at least two rounds to split the permutation into a top and
bottom half; it was skipped at one round.  At two rounds, five configurations
of $(r_0, \delta, \gamma)$---input difference byte, output bit mask, and
split---were probed over $200{,}000$ trials each.  The $5\sigma$ bias
threshold at this sample size is $0.011180$.  The highest observed correlation
was $|$corr$| = 0.003380$ ($r_0 = 1$, $\delta = (\text{byte}~0,\
\texttt{0x01})$, $\gamma_{\mathit{bit}} = 0$), less than a third of the
detection threshold.  All five configurations were below the threshold; no
differential-linear bias was detected.

\subsection{Test 5: Invariant subspace and S-box coset check}

\paragraph{Part A (dimension-1 S-box cosets).}
An exhaustive enumeration of all $(a, g)$ pairs satisfying
$S(a) \oplus S(a \oplus g) = g$ found $282$ hits across $141$ distinct
generators (out of $255$).  These are the candidate seeds of a dimension-1
invariant subspace trail.

\paragraph{Parts B and C (full-round survival).}
None of the $282$ hits survived a full round in the constant-base probe
(Part~B), nor in the random-base empirical probe (Part~C, $2{,}000$ trials per
sampled pair).  Both tests were run at one and two rounds with identical
conclusions: the S-box coset structure does not extend through a full round of
the permutation.

\paragraph{Part D (dimension-2 exhaustive check).}
All $32{,}385$ dimension-2 subspaces were checked exhaustively; no coset hit
was found.  The invariant trail therefore cannot be extended to higher dimension
via the S-box layer.

The $282$ dimension-1 hits are not unusual for an $8$-bit S-box and do not
constitute a weakness: the decisive test is whether the structure survives the
full round function, and it does not.

\section{Bit-level diffusion: SAC, completeness, and monobit}
\label{sec:sac}

The structural sanity checks below report \emph{mean} avalanche (a single input
flip changes $\approx0.5$ of the $2048$ output bits). A correct mean is
necessary but not sufficient for good diffusion: it is compatible with a
structured per-bit dependency in which an input bit deterministically flips some
output bit and never flips another. We therefore measured the full
strict-avalanche-criterion (SAC) matrix (Webster and Tavares, 1986),
\[
  \mathrm{SAC}[i][j] = \Pr\!\big[\text{output bit }j\text{ flips}\mid
  \text{input bit }i\text{ flipped}\big],\qquad\text{ideal }\tfrac12,
\]
over all $2048\times2048=4.19\times10^{6}$ (input,output) bit pairs, estimated
with $N=4096$ random base states per input bit at each round count from one to
eight (run against the AVX2 permutation; $2\times2048\times N$ permutation
evaluations per round). From the same counters we extract \emph{completeness}
(each cell strictly between $0$ and $N$: every input bit influences every output
bit, none deterministically) and a per-output-bit \emph{monobit} balance over
random inputs. At $N=4096$ the two-sided $5\sigma$ bound on a single proportion
is $|p-\tfrac12|>0.0391$.

\begin{table}[h]
\centering
\renewcommand{\arraystretch}{1.2}
\begin{tabular}{lccccc}
  \toprule
  \textbf{Round} & \textbf{SAC max} & \textbf{SAC mean} &
  \textbf{cells $>5\sigma$} & \textbf{dead / stuck} & \textbf{monobit max} \\
   & $|p{-}\tfrac12|$ & $|p{-}\tfrac12|$ & (exp.\ ${\approx}2.4$) & cells & $|f{-}\tfrac12|$ \\
  \midrule
  1 & 0.0415 & 0.0062 & 4 & 0 / 0 & 0.0276 \\
  2 & 0.0405 & 0.0062 & 2 & 0 / 0 & 0.0295 \\
  3 & 0.0422 & 0.0062 & 5 & 0 / 0 & 0.0305 \\
  4 & 0.0422 & 0.0062 & 2 & 0 / 0 & 0.0249 \\
  5 & 0.0413 & 0.0062 & 3 & 0 / 0 & 0.0264 \\
  6 & 0.0427 & 0.0062 & 1 & 0 / 0 & 0.0293 \\
  7 & 0.0435 & 0.0062 & 4 & 0 / 0 & 0.0271 \\
  8 & 0.0413 & 0.0062 & 4 & 0 / 0 & 0.0303 \\
  \bottomrule
\end{tabular}
\caption{SAC matrix, completeness, and monobit balance ($N=4096$ per input bit;
$5\sigma=0.0391$). The mean cell deviation $0.0062$ is exactly the value expected
of an ideal $\mathrm{Binomial}(N,\tfrac12)$ proportion ($\approx0.399\,\sigma$);
the $1$--$5$ cells per round exceeding $5\sigma$ track the $\approx2.4$ false
positives expected over $4.19\times10^{6}$ cells. Zero dead and zero stuck cells
at every round establish full bit-level completeness.}
\label{tab:sac}
\end{table}

At every round, including round~$1$, the SAC matrix is statistically
indistinguishable from ideal: the mean cell deviation equals the random-sampling
value and the few cells beyond $5\sigma$ match the expected false-positive count
over $4.19\times10^{6}$ cells. The completeness result is the decisive one---with
no dead and no stuck cells at any round, every one of the $2048$ input bits
provably influences every one of the $2048$ output bits, and none does so
deterministically. This is the per-bit form of the full-state first-round
diffusion that the mean-avalanche figure only asserts on average, and a design
with a correct mean but a structured dependency matrix would have failed it. No
output-bit monobit bias was detected at any round. The test reuses the
deterministic splitmix64 seeding convention of the other harnesses (seeded by
the round count) and is reproducible.

\section{Statistical randomness of the sponge keystream (NIST SP 800-22)}
\label{sec:nist}

The active-S-box bounds and the bit-level diffusion results above concern the
permutation directly. As a complementary mode-level check we ran the NIST
Statistical Test Suite (SP~800-22) on the keystream produced when the permutation
is used in the stream-cipher mode of the Krakken-Disk volume encryptor, where
$\mathit{ciphertext}_i = \mathit{plaintext}_i \oplus
\mathrm{Sponge}(\mathit{key}\,\|\,\mathit{nonce}\,\|\,\mathrm{LE64}(i))$.
Three independent $5500$\,MB volume containers were generated empty, so the
plaintext is the all-zero sector image and the resulting ciphertext is---up to
per-sector random nonces ($24$ B per $4096$ B sector) and a small fixed
header---the raw Krakken sponge keystream. The randomness test therefore reflects
the permutation output, not any plaintext structure.

Each file was assessed over $10$ sequences of $10^{6}$ bits. All fifteen STS test
families passed for all three files, each with a proportion of passing sequences
at or above the NIST minimum ($8/10$, or $5/6$ for the random-excursion tests)
and non-degenerate uniformity $P$-values: Frequency, Block Frequency, Cumulative
Sums, Runs, Longest Run, Binary Matrix Rank, Spectral (FFT), Non-Overlapping
Template ($148$ templates), Overlapping Template, Maurer's Universal, Approximate
Entropy, Random Excursions and Variant, Serial, and Linear Complexity. A custom
high-throughput monobit and runs test over a larger $10^{8}$-bit window of each
file returned unremarkable $P$-values (monobit $0.72$, $0.22$, $0.99$; runs
$0.53$, $0.31$, $0.53$ for the three files) and a byte-level Shannon entropy
above $7.999$ bits/byte in every case.

This is an empirical, mode-level result: like the distinguisher battery it is
necessary but not sufficient, excluding detectable first- and second-order
statistical bias in the keystream at the sampled scale rather than establishing a
property of the volume-encryption construction itself.

\section{Sponge-hash collision and output-distribution battery}
\label{sec:collision}

Complementing the keystream randomness tests, the sponge \emph{hash} ($256$-bit
digest, rate $160$ bytes) was probed for collision and distribution behaviour---a
family neither the active-S-box analysis nor the distinguisher battery
addresses---with three quick tests against the compiled hash.

\textbf{Birthday collisions.} Hashing $N=10^{6}$ random messages, truncating each
digest to $t$ bits, and counting colliding pairs reproduces the ideal
random-function expectation $\binom{N}{2}/2^{t}\approx N^{2}/2^{t+1}$
(Table~\ref{tab:birthday}). Observed counts match expectation closely wherever it
is statistically meaningful ($t\le34$: $1835$ vs $1863$, $495$ vs $466$, $122$ vs
$116$, $27$ vs $29$); at $t\ge36$ the expectation falls below $7$ and the ratios
are small-count Poisson noise. No truncation width shows a collision excess.

\begin{table}[h]
\centering
\renewcommand{\arraystretch}{1.15}
\begin{tabular}{rrrr}
  \toprule
  \textbf{$t$ (bits)} & \textbf{expected} & \textbf{observed} & \textbf{obs/exp} \\
  \midrule
  28 & 1862.6 & 1835 & 0.985 \\
  30 &  465.7 &  495 & 1.063 \\
  32 &  116.4 &  122 & 1.048 \\
  34 &   29.1 &   27 & 0.928 \\
  36 &    7.28 &    5 & 0.687 \\
  38 &    1.82 &    2 & 1.100 \\
  40 &    0.46 &    1 & 2.199 \\
  \bottomrule
\end{tabular}
\caption{Truncated-digest birthday collisions, $N=10^{6}$ messages, against the
ideal $N^{2}/2^{t+1}$.}
\label{tab:birthday}
\end{table}

\textbf{Output diffusion / near-collisions.} For $2\times10^{5}$ message pairs
differing in one input bit, the $256$-bit digest difference had mean Hamming
weight $127.98$ (ideal $128$), standard deviation $8.00$ (ideal $8$), minimum
$93$ and maximum $165$. The minimum far from $0$ excludes cheap one-bit
near-collisions.

\textbf{Digest byte uniformity.} Over $6.4\times10^{6}$ output bytes the
byte-value $\chi^{2}$ was $273.8$ on $255$ degrees of freedom (ideal
$255\pm22.6$; $z=+0.83$), consistent with a uniform output.

All three are empirical and sampling-bounded: they exclude gross collision or
distribution defects at the $\sim\!10^{6}$-message scale, not establish collision
resistance analytically.

\section{Remaining work}

The differential question that motivated this analysis is settled: the complete
round function has a proven minimum of $229$ active S-boxes at eight rounds, for
both differential and linear cryptanalysis. What remains, in priority order,
strengthens the analysis but is not expected to change the conclusion:

\begin{enumerate}\itemsep2pt
  \item \textbf{Independent bit-level linear model.} Replace the activity-level
        (transpose-coincident) linear model with an explicit bit-level mask
        model, for a check not sharing the differential model's structure.
  \item \textbf{Carry-aware ARX model.} Model PRESSURE's modular addition with
        the standard addition inequalities rather than activity-spreading.
  \item \textbf{Cluster / linear-hull estimates}, to complement the single-trail
        bounds.
  \item \textbf{Independent verification and external cryptanalysis.} Publish the
        specification, the MILP models, and the proven optima (each solver run is
        reproducible) for third-party scrutiny.
\end{enumerate}

\section*{Structural sanity checks}

Independently of the active-S-box bounds, the permutation was subjected to
empirical structural checks (not proofs): bijectivity on a sample
($2000/2000$ distinct outputs), absence of fixed points ($0$ among $503$ states
including the all-zero and all-one states), strict avalanche (single input-bit
flip changes a mean of $0.5004$ of the $2048$ output bits, minimum $0.47$), and a
sampled subspace-trail collapse test (random low-dimensional affine subspaces
show no rank collapse under the permutation). All checks pass. These rule out
several gross structural weaknesses but are not a substitute for an analytical
invariant-subspace / subspace-trail study, which remains future work.

\section*{Reproducibility}

The analysis tooling accompanying this note comprises: a scalar/AVX2 equivalence
harness (the two implementations are bit-identical on the permutation at all
round counts and on the hash across all tested input and output lengths); the
byte-lane MILP bound prover (differential and linear; SPN-core, SPN+XRBD, and
complete-round modes); the decomposed SPN-core prover; a single-lane propagation
probe; the trail-extraction verifier used to confirm the two-round result and
the branch-number property; and the structural sanity checker. All bounds were
obtained with the SCIP solver; each run is reproducible and terminates with a
certificate of optimality (or, where noted, a proven lower bound).

\end{document}
